% Shared module: what changes for nonlinear equations (short, clearly an outlook).
The results above rest on one structure: candidate equations form the commutative ring
$\ring=\C[\tau,\xi]$, solutions form vector spaces, and the exact sequence behind \cref{thm:factor} (a Chinese
remainder theorem up to the finite overlap $\deltaP$) turns a
factorization of the symbol into a splitting of the solution. For a nonlinear equation
$u_t=\sum_n\mathcal A_n(u)$, with each $\mathcal A_n$ a polynomial in the jet $(u,u_x,u_{xx},\dots)$, this structure is lost
in three ways.
\begin{enumerate}[label=(\roman*)]
  \item \emph{No ring of candidates.} Products of nonlinear operators are not operators of the same kind
  in any useful sense, and there is no unique factorization. One fixes instead a finite
  \emph{dictionary} $\mathcal D$ of jet monomials of bounded degree and order, with span $\mathcal V$, and asks
  which relations $F(u,u_x,\dots)=0$ in $\mathcal V$ the data satisfy. In practice each monomial (for instance
  $u_t$, $u_x$, $u_{xx}$, $uu_x$, $u^2$) is evaluated at the sample points, giving a matrix $\Theta$ with one column per
  monomial; a relation $F=\sum_mc_m\,m(u)$ is a vector $c$ with $\Theta c\approx0$, so the candidate equations are the
  near-null space of $\Theta$.
  \item \emph{No superposition.} Solutions do not form a vector space, so ``the component of $u$ along a
  factor'' has no meaning; the groups of \cref{thm:factor} do not exist.
  \item \emph{Identification is relative.} Identification can only hold relative to $\mathcal V$ and modulo the
  differential consequences of $F$ (its derivatives $\partial_xF,\partial_x^2F,\dots$ are also relations in
  any rich dictionary). The precise statement is an exact sequence relating the relations of $u$ in $\mathcal V$
  to those generated by $F$; identification holds exactly when a jet-level obstruction vanishes.
\end{enumerate}

What survives is the interaction term. Definition~\eqref{eq:eint-def} still makes sense with the
nonlinear single-term flows, and its short-time expansion is
\begin{equation}\label{eq:jet}
  \Eint(t)=\frac{t^2}{2}\sum_{n\ne m}\mathcal A_n'[u_0]\big(\mathcal A_m(u_0)\big)+O(t^3),
\end{equation}
where $\mathcal A_n'[u_0]$ is the Fr\'echet derivative of $\mathcal A_n$ at $u_0$. For Burgers' equation
$u_t=-uu_x+\nu u_{xx}$ the two terms are $\mathcal A_1(u)=-uu_x$ and $\mathcal A_2(u)=\nu u_{xx}$, with Fr\'echet
derivatives $\mathcal A_1'[u](v)=-(vu_x+uv_x)$ and $\mathcal A_2'[u](v)=\nu v_{xx}$. Then
\[
  \mathcal A_1'[u_0]\big(\mathcal A_2(u_0)\big)=-\nu\,(u_0''u_0'+u_0u_0'''),\qquad
  \mathcal A_2'[u_0]\big(\mathcal A_1(u_0)\big)=-\nu\,(3u_0'u_0''+u_0u_0'''),
\]
and half their sum gives
\[
  \Eint(t)=-\nu\,\big(u_0u_0'''+2u_0'u_0''\big)\,t^2+O(t^3).
\]
The expansion holds only while the solution is smooth. The flow of the single term $-uu_x$ is the inviscid
Burgers equation, whose characteristics cross: for $u_0=\sin x$ they are $x=\xi+t\sin\xi$, and
$\partial x/\partial\xi=1+t\cos\xi$ first vanishes at $\xi=\pi$, $t^*=1$. After this breaking time the single-term flow
has a jump, $\Eint$ has a jump in $x$ whose height does not vanish as $\nu\to0$, and the interaction can only be measured in
$L^1$, where for a single shock it is small, of order $\nu$ up to logarithmic factors (a rate we do not prove here).
Whether such data still identify the equation relative to a dictionary depends on how many shocks they contain
compared with the degree of the dictionary; we do not pursue this here.

\paragraph{Variable coefficients}
If the coefficients depend on $(t,x)$, the operators form the Weyl algebra $\C[t,x]\langle\partial_t,\partial_x\rangle$, which
is not commutative ($\partial_x\,x-x\,\partial_x=1$). Annihilators become left ideals, factorizations are no longer unique,
and \cref{thm:factor} has no general analogue. Worse, if the coefficients may be arbitrary smooth functions, one solution
identifies nothing: any $u$ without zeros satisfies the first-order equation $\big(\partial_x-u_x/u\big)u=0$, whose
coefficient is built from $u$ itself. Identification becomes possible only inside a fixed, finite-dimensional class of
coefficients. For the Schr\"odinger equation $iu_t=-u_{xx}+V(x)u$, for instance, the potential is read off directly as
$V(x)=(iu_t+u_{xx})/u$ wherever $u\neq0$. In the spectral picture of \cref{sec:alg-pencil}, slowly varying coefficients
would make the rates depend on position, $\lambda(\kappa;x)$ (\cref{rem:varcoef}).

\paragraph{A possible route} The grouping idea can be carried over by replacing ``modes on the same
irreducible component of $\Var(P)$'' with ``jet samples satisfying the same relation in $\mathcal V$''. Locally
in time, \eqref{eq:jet} shows that a nonlinear flow is a linear one (its Fr\'echet derivative) plus a
quadratic correction, so one can group linearized modes on short windows and test, through the size
of $\Eint$, whether the groups persist. We do not pursue this here: without a ring there is no
analogue of $\deltaP$, of the projectors~\eqref{eq:lagrange} or of $\gN$, and these are what make the
linear algorithm verifiable.

\begin{table}[tbp]
\centering\small
\caption{What carries over beyond scalar linear equations with constant coefficients. ``Groups'' means a splitting of
a solution into pieces attached to pieces of the equation.}\label{tab:beyond}
\begin{tabularx}{\linewidth}{@{}>{\raggedright\arraybackslash}p{2.3cm}>{\raggedright\arraybackslash}X>{\raggedright\arraybackslash}X>{\raggedright\arraybackslash}X@{}}
\toprule
Setting & Candidate equations & Groups & What survives\\
\midrule
linear, constant coefficients & commutative ring $\C[\tau,\xi]$; ideals; unique factorization & factors of the symbol, overlap $\deltaP$ (\cref{thm:factor}) & everything in this paper\\
variable coefficients & non-commutative Weyl algebra; left ideals; factorization not unique & no general splitting theorem & identification inside a finite-dimensional class of coefficients; local rates $\lambda(\kappa;x)$ (\cref{rem:varcoef})\\
systems (e.g.\ Dirac) & matrices over the ring; modules instead of ideals & sums of kernels need extra conditions and can fail on non-convex domains & identification criteria for square non-singular systems\\
nonlinear & finite dictionary of jet monomials; differential ideal & none (no superposition) & identification relative to the dictionary; the interaction term and its expansion~\eqref{eq:jet}\\
\bottomrule
\end{tabularx}
\end{table}
